% Section 3 -- drafted by Claude (2026-10-11), revised after checks (2026-10-11, 2026-10-21); V5 added 2026-10-21.
\section{Method}\label{sec:method}

My baseline (Codabench 954621) predicts $(4p_0 + k)/(4 + n)$ for a pair with $k$ successes in $n$ acquired labels, where $p_0$ is a prior success rate from the subject's metadata: a beta--binomial update with prior strength 4 (Appendix~\ref{app:model}). V1 to V5 build on one another; V5 is the final entry, and Table~\ref{tab:formative} gives the Codabench submission of each version.

Throughout, $\theta_s$ is subject $s$'s ability on the benchmark being predicted. Labels on other benchmarks are not used, and item difficulties are centered within each benchmark, so the benchmark's mean difficulty is absorbed into $\theta_s$. I write $\sigma$ for the logistic function and $\kappa(v) = \sqrt{1 + \pi v/8}$ for the factor of the probit approximation $\mathbb{E}[\sigma(x)] \approx \sigma(\mu/\kappa(v))$ for $x \sim \mathcal{N}(\mu, v)$ \citep{mackay1992evidence}.

\subsection{Ability prior and grid posterior (V1, V3)}
V1 gives every item of a pair the success probability $\sigma(\theta_s)$, with a normal prior $\theta_s \sim \mathcal{N}(m_s, s_\theta^2)$, $s_\theta = 1.5$, so the success rate is logit-normal. The prior mean is $m_s = \lambda\,(g(s) - \bar g)$, clipped to $[-6, 6]$, where $g(s)$ is the subject part of the baseline's metadata regression (release date, a ``thinking'' variant flag and the provider, plus a shrunken residual for subjects seen in training; Appendix~\ref{app:model}) and $\lambda$ is a scale. After $k$ successes in $n$ own labels,
\[
  p(\theta \mid \text{labels}) \propto \mathcal{N}(\theta; m_s, s_\theta^2)\, \sigma(\theta)^k\, \sigma(-\theta)^{n-k},
\]
which I evaluate on a fixed grid and summarize by the posterior mean of $\sigma(\theta)$, an expected-a-posteriori estimate \citep{bock1982eap}. V1 used $\lambda = 0.5$, although $\lambda = 1$ was better offline by 0.0016 (SE 0.0004): in a scenario analysis that is not in the repository, the half-strength prior had about a third of the full prior's worst-case loss when metadata do not carry over, with the same average. Two of its scenarios were modeled on the baseline's formative per-budget curve; the choice held without them. For V3 I set $\lambda = 0$; as no labels are passed at budget 0, every prediction there is exactly 0.5. This is the one model setting I chose directly from formative feedback, and offline it was worse (Section~\ref{sec:data}); the other use of feedback was the choice between V4 and V5 (Section~\ref{sec:results}). I tried one alternative to $s_\theta = 1.5$: at $\lambda = 1$, $s_\theta = 1.75$ scored $+0.0016$ (SE 0.0002) worse, and I kept 1.5.

\subsection{Item pooling (V2)}
In V2 items differ in difficulty. A Rasch model \citep{rasch1960} gives
\begin{equation}
  P(y_{si} = 1) = \sigma(\kappa_0\theta_s - z_i), \qquad z_i \sim \mathcal{N}(0, \tau^2), \qquad \kappa_0 = \kappa(\tau^2),
  \label{eq:rasch}
\end{equation}
with $\tau = 2$ (not tuned) and $\kappa_0 \approx 1.60$; a larger $z_i$ is a harder item. With $\phi_s = \kappa_0\theta_s$ this is the organizers' $\sigma(\phi_s - z_i)$ \citep{truong2025reeval}. The factor $\kappa_0$ keeps $\theta$ on V1's scale: by the probit approximation, averaging over an unknown difficulty gives back $\sigma(\theta_s)$, with an error of at most 0.0101 for $\theta_s \in [-6, 6]$.

A difficulty can only be learned from other subjects' labels on the same item, which the predictor sees when a draw puts several subjects on one benchmark. In my own formative draws few pairs shared a benchmark, so pooling seldom had partner labels to use. I estimate difficulties with a two-stage calibration \citep{ban2001online,truongkoyejo2026aims}: each partner label becomes an observation of the item's difficulty, given that partner's ability estimated from its \emph{other} labels, and the observations are combined with the prior on $z_j$ into a posterior mean $\hat z_j$ and variance $\omega_j^2$ (Appendix~\ref{app:model}). The target subject's labels on items with a known difficulty then enter its likelihood through $\hat z_j$ and $\omega_j^2$, its other labels enter as in V1, and the prediction for an item averages $\sigma(\kappa_0\theta - z)$ over the posteriors of $\theta$ and of the item's difficulty. Items carry no identifier, so I match them across subjects by their content and features. With no partner labels, V2 gives exactly V1's predictions.

\subsection{LLM difficulty ratings and the gate (V4)}
For V4 I also had Claude Opus 5.5 \citep{anthropic2026opus55} rate how hard each item is. The rater sees the item's content and features with a fixed rubric (Appendix~\ref{app:prompt}), is told not to solve the item, and returns one integer from 1 (almost every AI system released in 2023--2026 completes it) through 5 (about half do) to 10 (almost none do); a refusal or a value outside 1--10 counts as unrated. Offline I had it rate 1,049 public items, one per problem, which through shared content cover the 1,379 rated items of Table~\ref{tab:data}, at an estimated cost of \$12.08. Against Rasch difficulties fitted separately on each benchmark to all subjects' public responses (Appendix~\ref{app:model}), the Spearman correlations were 0.82 (366 rated problems), 0.28 (400 items), 0.44 (233) and 0.77 (50) on MathArena, Multi-SWE-bench, REAL and ResearchCodeBench. These benchmarks are public, and the rater may partly remember how hard their items are.

A rated item $j$ with rating $r_j$ gets the prior
\begin{equation}
  z_j \sim \mathcal{N}\bigl(\beta_t\,(r_j - \bar r),\ \rho_t^2\bigr),
  \nonumber
\end{equation}
which replaces the prior $\mathcal{N}(0, \tau^2)$ of Eq.~\eqref{eq:rasch} in the second stage and is used directly when the item has no partner labels; this is an explanatory item model in the sense of \citet{deboeck2004explanatory}. Here $\bar r$ is the mean rating of the rated items already labeled on the benchmark, by the target subject or by partners. Ratings thus enter only through their difference from that mean, and their absolute level is never used; until a rated item is labeled, and so always at budget 0, they have no effect. The slope $\beta$ and residual standard deviation $\rho$ come from a regression of the Rasch difficulties on the ratings, pooled within benchmarks: $\beta = 0.788$ and $\rho = 1.904$ for the submission, which uses all four rated benchmarks; the offline evaluation refits them without the held-out benchmark. A trust level $t \in [0, 1]$ interpolates between V3's prior and the fitted regression,
\begin{equation}
  \beta_t = t\beta, \qquad \rho_t^2 = \tau^2 - t^2 \max(\tau^2 - \rho^2, 0).
  \label{eq:trust}
\end{equation}
Because $\rho < \tau$ in every fit, $t = 0$ gives V3 and $t = 1$ gives $(\beta, \rho)$; because $\rho$ is close to $\tau$, $t$ acts almost only through the slope. I chose $t = 0.5$ ($\beta_t = 0.394$, $\rho_t = 1.976$) with a rule fixed before the run: it had the smallest worst-case regret over real, half-noisy and shuffled ratings (Section~\ref{sec:results}). A trust level's regret in a scenario is its Brier ALC minus that of the best trust level in the same scenario; half-noisy ratings keep each item's rating with probability 1/2 and otherwise take the rating of a random other item of the benchmark. With these values, after one solved label on an item rated 3, an item rated 8 gets 0.417 under V4 and 0.647 under V3.

Ratings that point the wrong way hurt. Offline, V4 without the gate was worse than V3 by 0.027 Brier ALC with the ratings reversed and by 0.0065 with them shuffled, and with the gate by 0.011 and 0.0048. The gate makes V4 predict as V3 whenever the benchmark has at least three labels on rated items, both outcomes occur among them, and the failed items have the lower mean rating. During scoring V4 requests ratings of the hidden items from the API with the same prompt; without a working key it predicts exactly as V3 (Section~\ref{sec:llm}).

\subsection{Guessed subject priors (V5)}
Ratings say nothing at budget 0, and on all three hosted draws scored with a metadata prior, the budget-0 predictions did worse than a constant 0.5 (Section~\ref{sec:results}). V5 instead asks Claude how strong the subject is. The request shows the subject's recorded fields, exactly as the platform passes them, and one item of the benchmark as an example, and asks what percentage of the benchmark's items the subject completes (Appendix~\ref{app:prompt}); it names no benchmark and contains no labels. During scoring the example is whichever item of the pair is predicted first. Before building V5 I checked, with a rule fixed in advance, whether Claude knows subject strength at all. With three example items per benchmark, the Pearson correlation between the logit of its guess and Rasch ability, computed within each benchmark and averaged over the four, was 0.818 (0.957 on MathArena, 0.767, 0.661 and 0.886 on the others). V1's metadata score reached 0.611 on the same measure. The single-example guesses that V5 uses reach 0.771.

The guess $q_s$ enters through a calibration. On the public pairs, an ordinary least-squares fit of the logit of each pair's success rate on $\operatorname{logit} q_s$, both clipped to $[0.01, 0.99]$, gives an intercept $a_g$, a slope $c_g$ and a residual standard deviation $s_g$. A subject with a guess then gets
\begin{equation}
  \theta_s \sim \mathcal{N}\bigl(t_g\,(a_g + c_g \operatorname{logit} q_s),\ (1 - t_g^2)\,s_\theta^2 + t_g^2 s_g^2\bigr),
  \label{eq:guessprior}
\end{equation}
with $q_s$ clipped to $[0.01, 0.99]$ and the mean to $[-6, 6]$ as in V1. This replaces V4's $\mathcal{N}(0, s_\theta^2)$ for the target subject and for every partner whose labels are pooled. The guess trust $t_g$ interpolates as in Eq.~\eqref{eq:trust}: $t_g = 0$ is V4 and $t_g = 1$ trusts the calibration fully. A subject without a guess keeps V4's prior. I chose $t_g = 0.75$ by the same minimax-regret rule as the rating trust, over real guesses, shuffled guesses and every guess moved by $\pm 1$ on the logit scale (Section~\ref{sec:results}). The submission's calibration uses all four rated benchmarks (204 pairs): $a_g = 0.040$, $c_g = 0.851$, $s_g = 0.811$, which gives the prior $\mathcal{N}(0.030 + 0.639 \operatorname{logit} q_s,\ 1.164^2)$. At budget 0 every item of a pair gets the prior mean of $\sigma(\theta_s)$: a guess of 10\% gives 0.25, 30\% gives 0.40 and 90\% gives 0.76, closer to 0.5 than the guesses themselves.
